\section{Preference Data: convertir «mejor» en etiquetas entrenables}

Una demonstration le indica al modelo cómo responder; la preference data, en cambio, pide a una persona o a un judge comparar varias respuestas. Añade grados de libertad como la task entropy, la longitud de la conversación, la prioridad entre valores y la rating scale, por lo que no se puede copiar sin cambios el proceso de SFT data collection.

\subsection{La interfaz de anotación de Helpfulness y Harmlessness}

Una preference sample suele contener varias model responses al mismo prompt, junto con una etiqueta chosen/rejected o rank/margin. El anotador no tiene que reescribir la respuesta, pero sí debe interpretar de forma estable qué significa «más helpful, más truthful o más harmless».

\lecturefigure{slide-23.jpg}{Interfaz de helpfulness y harmlessness para Human preference data}{Diapositivas oficiales, p.23}

\begin{knowledgebox}{Lectura de la figura: la label es el resultado de una comparación, no la verdad de un hecho}
La interfaz reduce el juicio humano a una pairwise choice o a una ordinal score. En preguntas factuales, la exactitud puede verificarse; en redacción, consejos y brainstorming, la preferencia depende del estilo, la longitud y los valores personales. Por eso la preference label es una protocol-dependent measurement.
\end{knowledgebox}

\lecturefigure{slide-24.jpg}{Las cinco desiderata de un preference dataset}{Diapositivas oficiales, p.24}

\begin{importantbox}{Las cinco perillas de la preference collection}
La task distribution determina qué se compara; la length distribution determina la carga de lectura y el contexto; single/multi-turn determina cuánto depende del historial; helpfulness/honesty/harmlessness determina la prioridad entre valores; la rating/ranking scale determina la label resolution. Un cambio en cualquiera de estas perillas puede modificar el reward model.
\end{importantbox}

\teachervoice{El ponente considera explícitamente la preference data como «otro problema de diseño de datos». Si los mismos prompts se convierten en conversaciones multi-turn, si se cambia la tie-break rule entre valores o si se cambia la rating scale, se obtienen señales de supervisión distintas.}

\subsection{Pilot: probar primero el flujo de anotación y después contratar a gran escala}

El equipo pidió primero a distintos vendors un pilot sobre 300 prompts de Self-Instruct, con la plantilla de Anthropic. El objetivo del pilot no era solo comparar proveedores, sino también detectar si la interface, la guideline o el model endpoint introducían un sesgo sistemático en las labels.

\lecturefigure{slide-25.jpg}{Pilot de preference collection con 300 prompts}{Diapositivas oficiales, p.25}

\begin{knowledgebox}{Lectura de la figura: el pilot es un measurement system test}
A la izquierda aparece la plantilla de anotación y a la derecha la tabla de resultados por vendor/task. Antes de ampliar la escala, conviene verificar si los distintos proveedores entienden la misma scale, si la tie rate es anómala y si algunas tasks son difíciles de juzgar. Si la measurement no es estable, agregar muestras solo repite el sesgo con mayor precisión.
\end{knowledgebox}

\lecturefigure{slide-26.jpg}{Diagnostics de task confusion/distribution en la preference data del pilot}{Diapositivas oficiales, p.26}

\begin{knowledgebox}{Lectura de la figura: la matriz muestra la cobertura y la confusión entre tareas}
La zona diagonal representa la intended category; las entradas fuera de la diagonal señalan inconsistencias en la task definition o en la prompt classification. Primero hay que ver qué categorías se confunden entre sí y después determinar si el label noise proviene del annotator, de la taxonomy o del prompt mismo; la matriz, por sí sola, no demuestra la calidad de las respuestas.
\end{knowledgebox}

\subsection{20K Dialogues, 80K Prompts y restricciones de Multi-turn}

La collection a gran escala tenía como meta unos 20K dialogues, con un total aproximado de 80K prompts. La distribución de tareas incluye a propósito distintos niveles de entropy, como generation, coding, math y QA; cada conversación completa tiene menos de 2,048 tokens y un promedio de unos cuatro turnos, para ajustarse al context de los modelos de ese momento y conservar follow-ups realistas.

\lecturefigure{slide-27.jpg}{Los 20K dialogues del preference dataset y la asignación de tareas}{Diapositivas oficiales, p.27}

\begin{knowledgebox}{Lectura de la figura: 20K y 80K describen unidades de conteo distintas}
20K es el número de dialogues y 80K, el número de prompts dentro de las conversaciones de varios turnos; con un promedio de cuatro turnos, ambos órdenes de magnitud se relacionan de forma natural. Todo reporte de entrenamiento o de costos debe indicar la unit; de lo contrario, una conversación de cuatro turnos se contará por error como una muestra o como cuatro muestras inconsistentes.
\end{knowledgebox}

\teachervoice{El equipo agregó tareas de baja entropía, como math, coding y closed QA, porque en ellas es más fácil juzgar si una respuesta es correcta; la creative generation se parece más a una conversación real, pero puede admitir muchas respuestas razonables, lo que reduce la concordancia entre las preferencias humanas.}

\begin{lstlisting}[caption={Estructura conceptual para convertir una multi-turn chosen/rejected conversation en un lote de DPO}]
def build_preference_example(dialogue, chosen, rejected, margin):
    return {
        "prompt": render_history(dialogue),
        "chosen": chosen,
        "rejected": rejected,
        "preference_margin": margin,
        "turn_count": len(dialogue),
    }
\end{lstlisting}

\lecturefigure{slide-28.jpg}{Longitud, número de turnos, valores y rating policy de la preference collection}{Diapositivas oficiales, p.28}

\begin{knowledgebox}{Lectura de la figura: las constraints sirven al modelo y a los anotadores a la vez}
El 2,048-token cap evita exceder el context del modelo y también limita el costo de lectura de la anotación; el promedio de cuatro turnos aporta follow-ups sin que la conversación completa se vuelva demasiado larga. Que el equipo priorizara helpfulness sobre harmlessness es una capability-stage choice y no debe interpretarse como un principio universal sobre la prioridad de la seguridad.
\end{knowledgebox}

\teachervoice{La preference collection no es una entrega única: cada semana el equipo recibía datos, hacía fine-tuning del modelo y entregaba al vendor un endpoint mejor para generar los pairs de la semana siguiente. La policy que enfrenta el annotator cambia continuamente, por lo que el dataset es un co-evolving process.}

\subsection{Chosen, Rejected y Preference Margin}

La subsección anterior describió la escala, el número de turnos y la prioridad entre valores de la collection; esta sección pasa al interior de un registro de entrenamiento y explica qué información aporta cada uno de winner, loser y margin. Al leer los ejemplos, primero hay que localizar el punto donde las dos respuestas realmente difieren y después ver si el margin corresponde a la magnitud de esa diferencia. Los ejemplos multi-turn mostrados en clase subrayan que tanto la chosen como la rejected pueden ser buenas; el margin expresa la intensidad de la preferencia, no una puntuación absoluta. Los pairs con margin grande sirven para aprender errores evidentes; los de margin pequeño contienen señales finas de style/value, pero también es más probable que incluyan annotator noise.

\lecturefigure{slide-29.jpg}{Multi-turn chosen/rejected preference examples y margin}{Diapositivas oficiales, p.29}

\begin{knowledgebox}{Lectura de la figura: primero el punto de divergencia, después el margin}
El ejemplo de la izquierda difiere claramente en la imitación de un personaje y en el manejo de una posible amenaza, y su margin es grande; el de la derecha es un texto de aniversario en el que ambas respuestas son utilizables y solo difieren ligeramente en detalles omitidos y en la redacción. La información de los Pairwise data proviene de «dónde difieren», no solo del texto del winner.
\end{knowledgebox}

\teachervoice{El equipo usó primero la scale de 1 a 8 de Anthropic y después la cambió a un esquema tipo Llama-2 de winner + margin de 1 a 4, porque este último explica con mayor facilidad «cuál es mejor y por cuánto». Un Scale change modifica la label distribution, por lo que también es una variable experimental.}

\subsection{Resumen de la sección}

La Preference data reduce la calidad de respuestas abiertas a etiquetas comparativas protocol-dependent. Una collection eficaz requiere diseñar en conjunto la task entropy, el multi-turn context, la tie-break rule entre valores, la rating scale y los endpoints iterativos.
